\section{Construct the arithmetic columns and the level datum}
\label{sec:arithmetic-constructors}

This section proves the two conditional constructors retained in RF10:
\texttt{chGood\_of\_host} and \texttt{exists\_levCh\_of\_levelPrime}.
The first starts with the actual arithmetic host, its prescribed top weights,
and an injective Boolean code. The second starts with the analytic
\texttt{LevelPrime} assertion and the prescribed positive quantities
$\mathrm{du}_u,D_*$ satisfying their logarithmic window. We construct the
columns and the complete level record, including their integrality and
root-of-unity conditions. The existence of those input hosts, good top
weights, and analytic primes is not asserted here.

Write $n=R.n$, $f=2n$, $r=k+m-1$ for the literal rank record. For a scale
$s$ put
\[
 x=s.x,\quad \epsilon=s.\mathrm{eps},\quad
 \mathcal L=x^\epsilon,\quad b=\exp(\mathcal L),\quad
 X'=\exp(x^{\epsilon/3}),\quad
 a=xb,\quad b'=\frac{xb}{X'},\quad w=\frac{b}{nX'}.
                                                               \tag{AC.1}
\]
The letter $b_j$ below is a column coordinate, whereas $b$ in (AC.1)
is a positive real scale. In the arithmetic host let
$E=H.E$, $s_d=H.\mathrm{sqd}$, and
$K=\mathbb Q(s_d)$, where $s_d^2=-d$.
The literal host hypotheses give a prime $d\equiv7\pmod {16}$,
an isomorphism
\[
 \operatorname{Gal}(E/\mathbb Q)
       \simeq\{\pm1\}^2\times A,
 \qquad |A|=n,
\]
and the second sign is the action on $s_d$.
These statements, rather than an arbitrary choice of a CM type, are
used in AC3.

\paragraph{AC1: correct all roots of unity with one bounded shift.}
Let $E$ be a number field, $F$ any field, and
$\varphi_i:E\hookrightarrow F$ a finite family of field embeddings.
Fix one index $i_0$ and arbitrary integers $a_i$. Put
$w_E=|\mu(E)|$, the number of roots of unity in $E$.
There exists an integer $q$ with $0\leq q<w_E$ such that
\[
       \prod_i\varphi_i(z)^{a_i+q\,1_{i=i_0}}=1
                  \qquad(z\in\mu(E)).                       \tag{AC.2}
\]
Here $w_E$ is the literal \texttt{Units.torsionOrder}.
For completeness, roots of unity are finite: an element of order $h$
has cyclotomic degree $\phi(h)\leq[E:\mathbb Q]$, while
$\phi(h)^2\geq h/2$ by the prime factor calculation in SZ10.
There are thus finitely many possible orders and finitely many roots
for each order. A finite subgroup of a field's multiplicative group is
cyclic. Indeed, combining elements of maximal prime-power orders gives
an element whose order is the exponent $e$ of the subgroup; all its
elements solve $X^e-1=0$, so its order is at most $e$, and it is cyclic.
Choose a generator $\zeta$ of $\mu(E)$.

The element
$U=\prod_i\varphi_i(\zeta)^{a_i}$ satisfies $U^{w_E}=1$.
The $w_E$ distinct powers of $\varphi_{i_0}(\zeta)$ are all the roots
of $X^{w_E}-1$ in $F$. Hence
$\varphi_{i_0}(\zeta)^q=U^{-1}$ for one $0\leq q<w_E$.
This proves (AC.2) at $\zeta$. Every $z\in\mu(E)$ is $\zeta^j$,
and the left side at $z$ is the $j$th power of its value at $\zeta$.
This proves the full assertion for every root of unity, including
negative exponents. It is the substance of \texttt{exists\_shift}.

\paragraph{AC2: the quadratic-subfield norm of every root is one.}
We prove the following elementary statement, also present as
\texttt{norm\_root\_of\_unity\_quadratic\_subfield} in the source:
if $d>3$ is prime, $s_d^2=-d$ in a number field $E$, and
$4\mid[E:\mathbb Q]$, then
\[
                  N_{E/K}(z)=1\qquad(z\in\mu(E)),
            \qquad K=\mathbb Q(s_d).                         \tag{AC.3}
\]
This will establish the needed host norm property rather than simply
reuse the field \texttt{HostDer.norm\_mu}.

First, $[K:\mathbb Q]=2$. If $y=a+c s_d\in K$ has rational square,
the coefficient of $s_d$ in $y^2$ gives $2ac=0$. Thus a rational number
$u$ which is a square in $K$ is either a rational square or $-d$ times
a rational square. It follows that
\[
             -1,-2,-3,2\quad\hbox{are not squares in }K.
                                                               \tag{AC.4}
\]
For $-j$ with $1\leq j\leq3$, the only possible second case would be
$j=d c^2$ with $c\in\mathbb Q$. The exponent of the prime $d$ on the
right is odd and on the left is zero, a contradiction. The first
case is excluded by negativity. For $2$, the second case is negative
and the first contradicts the elementary irrationality of $\sqrt2$.

The only root in $K$ of odd order is $1$. Otherwise a prime divisor
$p$ of that order gives a primitive $p$th root in $K$. Its degree
$p-1$ is at most two, so $p=3$. Such a root $v$ obeys
$v^2+v+1=0$, whence $(2v+1)^2=-3$, contradicting (AC.4).

Next let $y^{2^a}=1$. If $y^2=1$, then $y=\pm1$ and
$N_{E/K}(y)=1$, since $[E:K]=[E:\mathbb Q]/2$ is even.
If $y^2\ne1$, some power $i$ of $y$ has $i^2=-1$.
Put $K_1=K(i)$, which is quadratic over $K$ by (AC.4), and put
$v=N_{E/K_1}(y)$. The order of $v$ is a power of two.
It cannot exceed four: if it did, a power $u$ of $v$ would satisfy
$u^4=-1$, giving
\[
                         (u-u^3)^2=2.
\]
An element of $K_1$ whose square lies in $K$ has square either a
square of an element of $K$ or its negative, by the same coefficient
argument in the basis $1,i$. We would obtain a square root of $2$
or $-2$ in $K$, again contradicting (AC.4).
Thus $v\in\{1,-1,i,-i\}$. Each of these four elements has norm one
from $K_1$ to $K$, since the nonidentity $K$-automorphism sends
$i$ to $-i$. Norm transitivity proves $N_{E/K}(y)=1$.

Finally write an exponent killing $z$ as $2^a o$ with $o$ odd.
The result for $y=z^o$ gives $N_{E/K}(z)^o=1$.
The odd-order result in $K$ gives (AC.3).
This proof requires neither a class-number assertion nor an
unramifiedness or prime-distribution assertion.

\paragraph{AC3: identify the CM-type product and permit parallel shifts.}
For the literal host, the Galois isomorphism above and the Galois
degree formula give $[E:\mathbb Q]=4n$. Thus AC2 applies.
Let $\Psi\subset\{\pm1\}^2\times A$ be the kernel of the second
sign, and let $e_T:T\simeq\Psi$ be any finite reindexing.
For $g\in\Psi$, the automorphism $\sigma_g=H.\mathrm{gal}^{-1}(g)$
fixes $s_d$ and hence $K$. Conversely an automorphism fixing $K$
fixes $s_d\ne0$, so its second sign is $+1$.
Consequently $g\mapsto\sigma_g$ is a bijection from $\Psi$ to
$\operatorname{Gal}(E/K)$. The ordinary norm formula gives
\[
       \prod_{\tau\in T}\sigma_{e_T(\tau)}(z)=N_{E/K}(z)
                     \quad\hbox{inside }E.                  \tag{AC.5}
\]
Compose with the fixed host embedding $\varphi_0:E\hookrightarrow
\mathbb C$. Its compositions with these automorphisms are precisely
$H.\mathrm{emb}(e_T(\tau))$. By AC2,
\[
             \prod_{\tau\in T}H.\mathrm{emb}(e_T(\tau))(z)=1
                          \qquad(z\in\mu(E)).               \tag{AC.6}
\]

Write $\mathrm{MU}(a)$ for the assertion that this product with
integer exponents $a_\tau$ is one for every root $z$.
Equation (AC.6) proves $\mathrm{MU}(1)$; raising it to any integer
power proves $\mathrm{MU}(c)$ for every constant integer $c$.
Multiplying the two root identities proves
\[
                 \mathrm{MU}(a)\Longrightarrow
                      \mathrm{MU}(a+c).                     \tag{AC.7}
\]
If rational weights equal $a_\tau+c$ with these integer exponents,
they satisfy the literal $\mathrm{MUQ}$ definition: the required
integral lift is exactly $a+c$. This proves
\texttt{muq\_of\_parallel}, including its integrality assertion.
The condition is a product identity for all roots of unity; parity
of the sum alone would not suffice.

\paragraph{AC4: construct every field of a coded column.}
Let $D,h\in\mathbb N$ satisfy
\[
        D\geq1,\qquad h\geq(t+3)/2,
\]
and let $U_0,U_L,B\in\mathbb Z$ satisfy
\[
             4D(k-2)\leq U_0,\qquad U_L<0,
                         -B\leq U_L.                        \tag{AC.8}
\]
For a Boolean word $c=(c_0,\ldots,c_{k-3})$, put
\[
 \begin{split}
 v_i&=4D(k-2-i) &&(0\leq i\leq k-2),\\
 u_0&=U_0+\tfrac12,\qquad
 u_{k-1}=U_L+\tfrac12,\\
 u_{\ell+1}&=v_\ell-\tfrac12-D\,1_{c_\ell=\mathrm{true}}
                                     &&(0\leq\ell\leq k-3),\\
 b_j&=B+2hj,\qquad e_j=B+(2j+1)h &&(0\leq j<m).
 \end{split}                                                 \tag{AC.9}
\]
These are the four actual data fields of an $\mathrm{HCParam}(R)$.
We verify its nine proof fields.

For each interior index, $v_\ell-u_{\ell+1}$ is $1/2$ or
$D+1/2$, while
$u_{\ell+1}-v_{\ell+1}$ is $4D-1/2$ or $3D-1/2$;
all four numbers are positive. At the first endpoint,
$u_0>v_0$ by (AC.8); at the last, $u_{k-1}\leq-1/2<0=v_{k-2}$.
These prove both strict interlacing fields. Also
$u_{k-1}=U_L+1/2>-B=-b_0$.
The gaps $e_j-b_j$ and $b_{j+1}-e_j$ both equal $h$, so both
required lower bounds are valid. Each $b_j,e_j,v_i$ is integral,
and each $u_i$ is an integer plus $1/2$.
This proves all fields, including the strict bottom inequality.
It proves \texttt{exists\_hcParam\_of\_code} by an explicit formula.

\paragraph{AC5: uniform numerical prerequisites for the column formula.}
Fix $0<\epsilon<1/52$, $c_D\geq0$, and $c_h\geq402$.
There is $x_0$, depending only on these three fixed numbers, such that
for every scale with this $\epsilon$, every $x\geq x_0$, and every
actual $R$ satisfying $\mathrm{RankVal}(s,R)$, set
\[
                  D=\lceil x^{c_D}\rceil_+,
             h=\lceil x^{c_h}\rceil_+,
                    h_0=2h.                                \tag{AC.10}
\]
Then
\[
 D\geq1,\quad (t+3)/2\leq h,\quad 4kD\leq a,
       \quad k\leq b'/x,\qquad b',b,w>0.                    \tag{AC.11}
\]
Indeed, use RF12 to obtain $k\leq x^5$ and $t\leq x^{401}$,
and require $x\geq8$. Since $D\leq2x^{c_D}$,
\[
       4kD\leq8x^{c_D+5}\leq x^{c_D+6}.
\]
RF11 lets us require $x^{c_D+6}X'\leq b$.
Since $X'\geq1$ and $a=xb$, this proves $4kD\leq a$.
It also proves $k\leq x^5\leq b/X'=b'/x$.
The estimate $h\geq x^{402}\geq(t+3)/2$ follows from
$t\leq x^{401}$ and $x\geq8$.
All positivity assertions follow from (AC.1) and $n\geq1$.
The threshold precedes the choices of $s,R$, the host, the code,
and the top weights. This is the exact uniformity of
\texttt{hc\_sizes}.

\paragraph{AC6: construct the whole column family with exact root congruences.}
Retain AC5, supply an actual host $H$ with
$\mathrm{HostVal}(s,C_8,H)$, and supply integers $a_T(t),b_T(t)$
on the host's top set satisfying
\[
                         a\leq a_T(t),\qquad b'\leq b_T(t).
                                                               \tag{AC.12}
\]
Let $e_T:T\simeq\Psi$ be any finite reindexing. We construct
\emph{simultaneously for every Boolean word} a family
$p_{c,\tau}\in\mathrm{HCParam}(R)$.

The top places are $((1,1),\iota(t))$, while thick places are
$((-1,1),\iota(s\,a_0+H.u_0))$ for $s\in H.T_2$.
Both lie in $\Psi$ and their first signs make them disjoint.
The inequality $|H.T_2|\geq x^{-\epsilon}n>0$ in
$\mathrm{HostVal}$ makes the thick set nonempty.
Choose a thick place $\tau_0\in T$.
For top places the index $t$ is unique, because $\iota$ is injective.
Define the integer base
\[
 B^{\mathrm{raw}}_\tau=
 \begin{cases}
 b_T(t)+\lfloor b'/x\rfloor,&e_T(\tau)=\tau_1(t),\\
 \lceil b\rceil,&e_T(\tau)\text{ is thick},\\
 \lceil w\rceil,&e_T(\tau)\text{ is thin}.
 \end{cases}                                                 \tag{AC.13}
\]
AC1 supplies $0\leq q<w_E$ such that
\[
             B_\tau=B^{\mathrm{raw}}_\tau+q\,1_{\tau=\tau_0}
                          \quad\hbox{satisfies }\mathrm{MU}(B).
                                                               \tag{AC.14}
\]
The correction changes no top or thin coordinate. On every thick
place it gives exactly
\[
                 \lceil b\rceil\leq B_\tau
                            <\lceil b\rceil+w_E.             \tag{AC.15}
\]
Thus the source's strict torsion bound is preserved.

Write $k=2\kappa$, which is possible because $R.k$ is even.
Choose endpoint integers
\[
 (U_0(\tau),U_L(\tau))=
 \begin{cases}
 (a_T(t)+\kappa-1,-b_T(t)-\kappa),
                                  &e_T(\tau)=\tau_1(t),\\
 (4D(k-2)+4kD,-1),                &\text{otherwise}.
 \end{cases}                                                 \tag{AC.16}
\]
At a top place, $a_T\geq4kD$ gives the first condition in (AC.8).
Also $b_T>0$ and
$\lfloor b'/x\rfloor\geq k\geq\kappa$, so $U_L<0$ and
$-B\leq-b_T-\kappa=U_L$.
Off the top set, $B\geq1$ by positivity of $b,w$ and $q\geq0$;
(AC.8) follows directly. Apply AC4 using these same $D,h,B,U_0,U_L$
for every word.

The resulting family has
\[
 \begin{split}
 v_i(\tau)&=4D(k-2-i),\\
 b_j(\tau)&=B_\tau+jh_0,\qquad
 e_j(\tau)=B_\tau+(j+\tfrac12)h_0,\\
 u_{\ell+1}(\tau)&=v_\ell(\tau)-\tfrac12
                         -D\,1_{c_\ell=\mathrm{true}}.
 \end{split}                                                 \tag{AC.17}
\]
Since $h_0=2h$, the weight $e_j$ is the integral weight $B+(2j+1)h$.
Equations (AC.14) and (AC.7) prove $\mathrm{MUQ}(e_j)$ for every $j$.
This supplies the exact root condition in $\mathrm{HCVal}$.
All its other fields are (AC.13), (AC.15), (AC.16), and (AC.17).
In particular, off the top set the endpoints are
$4D(k-2)+4kD+1/2$ and $-1/2$, and on a top place they are
\[
         u_0=a_T(t)+(k-1)/2,\qquad
         u_{k-1}=-b_T(t)-(k-1)/2.                            \tag{AC.18}
\]
The functions $v,b,e$ and both endpoint functions are independent
of the word. This proves every conclusion of
\texttt{exists\_hc\_family}, including the five compatibility
identities between any two words and their prescribed top values.
No unproved global simultaneous choice of these column families
remains under the stated inputs.

\paragraph{AC7: the literal \texttt{chGood\_of\_host} constructor.}
Fix $\epsilon,C_5,c_6,C_b,c_D,c_h,C$ as in that theorem:
$0<\epsilon<1/52$, $\mathrm{Std}(C)$, $c_D\geq0$, $c_h\geq402$.
Then a uniform threshold works for all the following supplied data:
an actual $R,H$, $\mathrm{RankVal}(s,R)$,
$H.CT_2=C.CT_2$, the literal
\[
 \mathrm{HostGood}(s,C.P.C_{\mathrm{pde}},C_5,c_6,C_b,
                                      H,q_0(H)),
                                                               \tag{AC.19}
\]
prescribed weights $a_T,b_T$ with $\mathrm{WinOK}$, a code
$c:\mathrm{Fin}(M)\hookrightarrow\{\mathrm{false},\mathrm{true}\}^{k-2}$,
and
\[
 M=\left\lceil3\Gamma^2x^{C.P.C_F}\right\rceil_+,
       \qquad \Gamma=\exp(C.C_\Gamma m^5f\mathcal L).
                                                               \tag{AC.20}
\]
An arbitrary natural $u$ may also be prescribed as the host's
unipotent translate.

Put $N_c=H.Q_S H.q H.d$ and $C_8=C.P.C_{\mathrm{pde}}\geq1$.
The pair-row part of (AC.19) gives
\[
              \log N_c\leq2x^{\epsilon/5}(\log x)^2.
\]
For every $z\leq N_c^{C_8}$, require the uniform elementary
inequality
$2C_8(\log x)^2< x^{2\epsilon/15}$.
Then $z<X'$. The host bounds therefore give
\[
                q_0(H),p_{\mathrm{de}},p'_{\mathrm{de}}<X'.
                                                               \tag{AC.21}
\]
On the other hand, each $g_i\geq(kb)^8\geq b\geq X'$.
Thus $q_0(H)\ne g_i$, with no need to assume this extra inequality.

The host's level-prime set is exactly
$\{g_0,g_1,p_{\mathrm{de}},p'_{\mathrm{de}},q_0(H)\}$ by
(AC.19). Every member is prime and congruent to $1$ modulo $N_c$:
for the two $g_i$ this is $\mathrm{HostArith}$; for the two
$p_{\mathrm{de}}$ it follows by reducing their stated congruence
modulo $\operatorname{lcm}(15,N_c)$; for $q_0(H)$ it is already
in (AC.19). Here $N_c>1$.
The same input supplies $H.cc_K=\mathcal O_K$ and the disjointness
of the two literal $\mathrm{BB}$ sets.

There is also a stronger congruence in \texttt{ChOK.small}.
The source defines $r_0(H)$ to be the least prime not dividing
$N_c(p_{\mathrm{de}}-1)(p'_{\mathrm{de}}-1)N(cc_K)$ and defines
$q_0(H)$ as the least prime $p$ with
\[
              p\equiv1\pmod {r_0(H)N_c},
                           \qquad p\nmid N(cc_K).
\]
Because $q_0(H)$ is already prime by (AC.19), this defining set is
nonempty: its natural infimum would be zero if it were empty.
The least element of a nonempty set of natural numbers belongs to
that set. Thus the stronger congruence follows directly from the
literal definition. Since $cc_K=\mathcal O_K$ has norm one,
$r_0(H)$ is exactly \texttt{leastOff}
$\bigl(N_c(p_{\mathrm{de}}-1)(p'_{\mathrm{de}}-1)\bigr)$.
This argument does not add a new prime-existence hypothesis.

Replace only $H.u_0$ by $R.\mathrm{ofNat}(u)$.
The definitions of $\mathrm{HostGood}$ and $\mathrm{WinOK}$ do not
read that field, so (AC.19) and all prescribed weights persist.
The thick set can change; construct the AC6 family after this
replacement, so its thick-place correction uses the correct new set.
The lower halves of $\mathrm{WinOK}$ give (AC.12).
Since the positive expression in (AC.20) has positive ceiling,
$M>0$. Choose the base word $0\in\mathrm{Fin}(M)$ and put
\[
                         \mathrm{hcW}(v,\tau)=p_{c(v),\tau}.
\]
Together with the supplied $R$, the translated $H$, and the preceding
level and $\mathrm{BB}$ proofs, these data fill the actual
\texttt{PreCh} record. Its five word-compatibility fields are exactly
those proved in AC6.

For precision, the sixteen fields of $\mathrm{ChOK}$ are supplied as
follows:
\[
\begin{array}{c|l}
\text{fields}&\text{proof or supplied value}\\ \hline
\mathrm{rank},\mathrm{ct2}&\mathrm{RankVal},\ H.CT_2=C.CT_2\\
\mathrm{val},\mathrm{arith},\mathrm{thick},\mathrm{dlog}
      &\text{the corresponding clauses of (AC.19)}\\
\mathrm{gnorm},\mathrm{pde}&\text{the integral norm and congruence clauses of (AC.19)}\\
\mathrm{small}&\text{level membership, (AC.21), and the stronger congruence above}\\
\mathrm{lev}&\text{the five-element level set and (AC.21)}\\
\mathrm{win},\mathrm{card},\mathrm{inj}
      &\mathrm{WinOK},\ \text{(AC.20)},\ \text{the supplied injective code}\\
\mathrm{hc},\mathrm{code},\mathrm{top}&\text{AC6, (AC.17), and (AC.18)}
\end{array}                                                   \tag{AC.22}
\]
In the \texttt{small} row, inequality with the two $p_{\mathrm{de}}$
is a clause of (AC.19), while inequality with the $g_i$ was proved
above. The retained $\mathrm{HostGood}$ is the second field of
$\mathrm{ChGood}$. This proves the literal endpoint with its
prescribed translate, weights and code. Its threshold depends only
on the fixed displayed constants, and is chosen before $R,H,a_T,b_T,M,c,u$.

\paragraph{AC8: turn an integral norm-form solution into the required element.}
For $d$ prime with $d\equiv3\pmod4$, write $s_d^2=-d$ and let
conjugation send $s_d$ to $-s_d$. The literal assertion
$\mathrm{NormOne}(d,M,p)$ means that there are $A,B\in\mathbb Z$
with
\[
 A^2+dB^2=4p,\qquad M\mid B,\qquad 2M\mid A-2-B.
                                                               \tag{AC.23}
\]
The pair-based and two-integer source definitions are identical.
Write $B=M u$, $A=2+Mu+2Mv$, and put $e=(d+1)/4\in\mathbb Z$.
The element $\omega=(1+s_d)/2$ is integral because
$\omega^2-\omega+e=0$. Define
\[
       \pi=1+M(v+u\omega)=\frac{A+B s_d}{2}\in\mathcal O_E.
                                                               \tag{AC.24}
\]
It lies in $K$, satisfies $\pi\overline\pi=p$, and
$\pi-1\in M\mathcal O_E$. Moreover expansion gives
\[
     p-1=M\left(u+2v+M(eu^2+uv+v^2)\right).
                                                               \tag{AC.25}
\]
Thus $p\equiv1\pmod M$. This proves every conclusion of
\texttt{normOne\_host}, with the actual element (AC.24), not just
an abstract splitting assertion.

The converse used in \texttt{normOne\_iff\_of\_sqrt} can also be
proved directly. Suppose $M>0$, $\pi\in\mathcal O_E\cap K$,
$\pi\overline\pi=p$, and $\pi-1=M y$ with $y\in\mathcal O_E$.
Then $y\in K$, so write $y=a+c s_d$ with $a,c\in\mathbb Q$.
Both its trace $2a$ and norm $a^2+dc^2$ are rational algebraic
integers and hence integers, say $t$ and $v_0$.
The identity $d(2c)^2=4v_0-t^2\in\mathbb Z$ implies
$2c\in\mathbb Z$: in lowest terms its denominator squared divides
the squarefree integer $d$. Write $2c=u$.
Reducing $t^2+du^2=4v_0$ modulo four and using $d\equiv3\pmod4$
shows that $t,u$ have the same parity. Set
$A=2+Mt$ and $B=Mu$. Then (AC.23) follows by taking the norm of
$\pi=1+My$, and $2M\mid A-2-B$ follows from the parity.
For $M=0$, both sides of the equivalence force $\pi=1$ and $p=1$;
use $A=2,B=0$. Finally, replacing $M$ by any divisor $M'$ preserves
(AC.23), so the divisor version of the congruence follows as well.

\paragraph{AC9: construct the full level record from its prime solutions.}
Supply an actual $P:\mathrm{PreCh}$, positive integers $\mathrm{du}_v$,
positive $D_*$, and $\Gamma>1$. Write
\[
 U_v=\left(\frac{D_*}{\mathrm{du}_v}\right)^{1/(fk)}
                         \Gamma^{-2/(fk)},\qquad
 M_c=N_c(g_0g_1)^2.                                         \tag{AC.26}
\]
Suppose for every word $v$ we have a prime $T_v$ and an integral
$\pi_v\in K$ with
\[
 \begin{gathered}
 \pi_v\overline{\pi_v}=T_v,\quad
 \pi_v-1\in M_c\mathcal O_E,\quad
 T_v\ne p_{\mathrm{de}},p'_{\mathrm{de}},\\
       U_v\leq T_v\leq(1+1/(fk))U_v.
 \end{gathered}                                              \tag{AC.27}
\]
These data construct the actual $\mathrm{LevCh}(P)$.
The sole additional arithmetic field to check is
$T_v\equiv1\pmod {N_c}$.
Since $N_c\mid M_c$, write $\pi_v=1+N_c\alpha$ with
$\alpha\in\mathcal O_E$. Then
\[
 \frac{T_v-1}{N_c}
       =\alpha+\overline\alpha+N_c\alpha\overline\alpha
\]
is both rational and an algebraic integer, so it is an integer.
Since $N_c>1$, the literal remainder is $1$.

Set the five data fields $T,\pi T,\mathrm{duStar},Dstar,GamP$
to $T_v,\pi_v,\mathrm{du}_v,D_*,\Gamma$, respectively.
The ten proof fields are: primality, the congruence just proved,
$\pi T\in K$, its norm identity, its congruence modulo $M_c$,
avoidance of the two specified primes, positivity of $\mathrm{duStar}$,
positivity of $Dstar$, $1<GamP$, and the two-sided window.
They are exactly (AC.27) and the stipulated positive values.
No pairwise distinctness of the $T_v$ is required by this record.
This proves \texttt{LevelChoice.exists\_levCh}.

The associated host has level set
$H.\mathrm{levelPrimes}\cup\{T_v\}$. Its old membership fields persist.
Each new prime is coprime to $Q_S$, since it is $1$ modulo the
multiple $N_c$ of $Q_S$. The $cc_K$ coprimality fields follow from
$cc_K=\mathcal O_K$. Thus this is also the literal
\texttt{LevCh.host} used by the later joins, with the same field,
rank, frame labels, and column data.

\paragraph{AC10: use exactly the retained analytic level-prime assertion.}
Assume the literal $\mathrm{LevelPrime}$ statement. For every fixed
$0<\epsilon<1/52$, $c_s>0$, and standard $C$, there is a threshold
$x_0$ with the following property, uniformly over \emph{all} $c_D,c_h$
and all subsequently supplied $s,P,q_1,a_T,b_T,c$.
Suppose $s.\mathrm{eps}=\epsilon$, $x\geq x_0$, and
$\mathrm{ChOK}(s,C,c_D,c_h,P,q_1,a_T,b_T,c)$.
Supply positive $\mathrm{du}_v,D_*$ with
\[
       c_s m\mathcal L\leq\log U_v\quad\text{for every }v,
       \qquad
       \Gamma=\exp(C.C_\Gamma m^5 f\mathcal L),               \tag{AC.28}
\]
where $U_v$ is (AC.26) for this $\Gamma$.
We construct $L:\mathrm{LevCh}(P)$ with the exact equalities
\[
              L.\mathrm{duStar}=\mathrm{du},\quad
              L.Dstar=D_*,\quad L.GamP=\Gamma.              \tag{AC.29}
\]

Here is the precise analytic input and how its hypotheses are met.
$\mathrm{LevelPrime}$ supplies, once its threshold depending on
$\epsilon,c_s$ is met, a prime $T$ in
$[U,(1+1/(fk))U]$ and
$\mathrm{NormOne}(d,N_c(g_0g_1)^2,T)$, provided
\[
 \begin{gathered}
 \mathrm{RankVal}(s,R),\quad d\text{ prime},\quad d\equiv3\pmod4,
       \quad d\geq100,\quad 0<N_c,\quad d\mid N_c,\\
 \log N_c\leq3x^{\epsilon/5}(\log x)^2,\qquad
 1\leq g_i\leq b^9,\quad U>0,\quad c_s m\mathcal L\leq\log U.
 \end{gathered}                                              \tag{AC.30}
\]
The conditions on $d$, $N_c$, and $\log N_c$ are the literal
$\mathrm{ModRow}(s,N_c,d)$; $\mathrm{RankVal}$ is a separate input.
The host supplies primality and $d\equiv7\pmod {16}$.
Its bound $4tr+1<d$, together with $r\geq6,t\geq11$, gives
$d>265>100$. The other two elementary conductor conditions are
immediate from $N_c=Q_Sqd$. The pair row inside $\mathrm{ChOK.val}$
has $1\leq Y\leq x$ and
$\log N_c\leq2Y^{\epsilon/5}(\log Y)^2$; monotonicity gives
the required stronger bound with coefficient two, hence (AC.30).

Require the RF11--RF12 thresholds, $x\geq2$, and
$x\geq(1/c_s)^4$. Then
\[
 k\leq x^5,\quad x^{41}\leq b,\quad m\geq x^{1/4},
        \qquad 2k^8\leq2x^{40}\leq x^{41}\leq b.
\]
The host's upper bound $g_i\leq2(kb)^8$ now gives $g_i\leq b^9$.
Its primality gives $g_i\geq1$.
The same thresholds imply $c_s m\geq1$, and therefore
\[
              X'\leq\exp(c_s m\mathcal L)\leq U_v.
                                                               \tag{AC.31}
\]
All $U_v$ are positive; $\Gamma>1$ because $C.C_\Gamma\geq1$
and $m,f,\mathcal L$ are positive.
For $p_{\mathrm{de}}$ and $p'_{\mathrm{de}}$, the field
$\mathrm{ChOK.lev}$ gives either equality with some $g_i$ or a
strict bound by $X'$. The equality is excluded by the literal host
axiom $p_{\mathrm{de}},p'_{\mathrm{de}}\ne g_i$.
Thus both specified primes are less than $X'$.

Apply $\mathrm{LevelPrime}$ for each word with $U=U_v$.
It supplies all the prime windows and integer norm solutions at
one common threshold: its quantifier over $U$ occurs after its
threshold, so neither the word nor $\mathrm{du},D_*$ changes $x_0$.
Use AC8 to produce $\pi_v$ with congruence modulo $M_c$.
Equation (AC.31) gives $T_v\geq U_v\geq X'$, excluding both
$p_{\mathrm{de}}$ primes. AC9 constructs the desired record and
(AC.29). Taking the maximum of these finitely many thresholds
proves exactly \texttt{exists\_levCh\_of\_levelPrime}.

Consequently RF10's two arithmetic constructor arrows are now proved.
They remain conditional on the same actual inputs: the supplied
$\mathrm{HostGood}$, top windows and injective code for AC7; and
$\mathrm{LevelPrime}$ plus the positive discriminant data and their
logarithmic window for AC10. The construction of recursive global
forms, the joint $\mathrm{RecExP}$ data, and the independent analytic,
spectral and pairing branches are not consequences of these arrows.
